\chapter{Actualizaciones}

Un release tiene partes que se actualizan por caminos distintos:

\begin{center}
\begin{tabularx}{\linewidth}{@{}l l X@{}}
\toprule
\textbf{Parte} & \textbf{Dónde vive} & \textbf{Cómo se actualiza} \\
\midrule
Bitstream de la FPGA & flash de configuración de la FPGA & JTAG (FlashPro), con acceso físico \\
Driver y servicio & tarjeta microSD & por red, con los instaladores \\
Librería cliente & PC del usuario & \cmd{pip install} del nuevo wheel \\
\bottomrule
\end{tabularx}
\end{center}

\begin{advertencia}[Limitación de la placa]
El PolarFire SoC permite actualizar el bitstream en campo (\emph{Auto Update}
e \emph{IAP}), pero ambos requieren una memoria flash SPI conectada al System
Controller del chip, que la Discovery Kit no tiene. En esta placa el
bitstream sólo se actualiza por JTAG. Una placa de producto debería incluir
esa memoria.
\end{advertencia}

\section{Compatibilidad}

\begin{itemize}
  \item \cmd{/v1/device} informa las versiones de las tres partes de la placa;
    \cmd{bitstream} identifica el release y el commit exacto.
  \item Dentro del protocolo v1, un cliente de cualquier versión funciona con
    un dispositivo de cualquier versión: el protocolo sólo crece por adición.
\end{itemize}

\section{Integridad de los archivos}

Cada release publica un archivo \cmd{SHA256SUMS}. Antes de instalar:

\begin{lstlisting}[language=shell]
sha256sum -c SHA256SUMS
\end{lstlisting}
